%% GENERATED FILE -- do not hand-edit.
%% SOURCE: experiments/database/scripts/build_candidate_datasets_table.py
\begin{landscape}
\begingroup
\footnotesize
\setlength{\tabcolsep}{4pt}
\renewcommand{\arraystretch}{1.15}
\begin{longtable}{@{}r@{\hspace{4pt}} L{86mm} L{92mm} L{62mm}@{}}
\caption[]{Sources for Table~\ref{tab:candidates}, in the same order. \emph{Why it is
ordered here} is the band reason for a perturb-seq or molecular-layers row, and the tier
and measurement rule for a scale row; for the first 60 rows what the row buys is the \emph{Why} column of Table~\ref{tab:final}. \emph{Data} is where the per-record values live; entries marked
unconfirmed were not fetched live and must be checked before a loader is written. Every
link is clickable.}
\label{tab:sources-all}\\
\toprule
\textbf{\#} & \textbf{Dataset, citation and link} & \textbf{Why it is ordered here} & \textbf{Data} \\
\midrule
\endfirsthead
\multicolumn{4}{@{}l}{\footnotesize\emph{Table~\ref{tab:sources-all}, continued}}\\
\toprule
\textbf{\#} & \textbf{Dataset, citation and link} & \textbf{Why it is ordered here} & \textbf{Data} \\
\midrule
\endhead
\bottomrule
\endfoot
1 & \textbf{Muenzner 2024 (natural-isolate proteome)}\newline Muenzner J, et al., Ralser M. Nature 2024;630:149-157.\newline \href{https://doi.org/10.1038/s41586-024-07442-9}{doi:10.1038/\allowbreak s41586-024-07442-9} & Next build, first: the proteome on the sequenced 1,011-isolate panel whose transcriptomes the supported Caudal 2024 already carries, so protein and mRNA sit on one genotype axis strain by strain. The natural-variation counterpart of the Kemmeren 2014 and Messner 2023 pair that Figure 3 rests on. & MassIVE MSV000090435 (natural collection); PRIDE PXD044526 (disomes), PXD048219 (turnover) \\
\addlinespace[5pt]
2 & \textbf{Peter 2018 (1,011 isolate genomes + phenome)}\newline Peter J, De Chiara M, Friedrich A, et al., Schacherer J. Nature 2018;556:339-344.\newline \href{https://doi.org/10.1038/s41586-018-0030-5}{doi:10.1038/\allowbreak s41586-018-0030-5} & Next build, second: the 36-condition growth phenome on the same panel, and the loss-of-function (nonsense plus SIFT) and homozygous-frameshift matrices already mirrored from the host. Those matrices define a natural loss-of-function perturbation, so the model sees a gene broken two ways, by engineered deletion and by a natural allele; the ontology has no type for natural gene loss yet and the core-loss copy-number variant stands in. & ENA/\allowbreak SRA; phenotype tables via the paper and France Genomique; assemblies, pangenome, copy-number, loss-of-function and frameshift matrices mirrored from 1002genomes.u-strasbg.fr \\
\addlinespace[5pt]
3 & \textbf{De Chiara 2022 (domestication and the life cycle)}\newline De Chiara M, Barre BP, Persson K, et al., Liti G. Nat Ecol Evol 2022;6:448-460.\newline \href{https://doi.org/10.1038/s41559-022-01671-9}{doi:10.1038/\allowbreak s41559-022-01671-9} & Next build, third: sporulation, germination and other life-cycle traits on the same panel, the natural-variation counterpart of the supported CalMorph morphology on deletions, so a cell-level phenotype is measured on both genotype axes. & Nat Ecol Evol SI (unconfirmed) \\
\addlinespace[5pt]
4 & \textbf{Boocock 2025 (single-cell eQTL mapping)}\newline Boocock J, Alexander N, Alamo Tapia L, Walter-McNeill L, Patel SP, Munugala C, Bloom JS, Kruglyak L. eLife 2025;13:e95566.\newline \href{https://doi.org/10.7554/eLife.95566}{doi:10.7554/\allowbreak eLife.95566} & Genotype recovered per cell from the cell's own reads, crossed with a per-cell transcriptome. The largest existing yeast dataset in the quadrant the campaign targets. & mirrored as boocockSinglecellEQTLMapping2025; GEO (unconfirmed) \\
\addlinespace[5pt]
5 & \textbf{Hale 2024 (CRISPRi x natural variation)}\newline Hale JJ, Matsui T, Goldstein I, et al., Kruglyak L. Nat Commun 2024;15:4234.\newline \href{https://doi.org/10.1038/s41467-024-48626-1}{doi:10.1038/\allowbreak s41467-024-48626-1} & The one released yeast library that crosses a CRISPR interference guide set with a sequenced genetic background, which is the input axis a campaign has to plan against. & SRA PRJNA986287 + Figshare \\
\addlinespace[5pt]
6 & \textbf{Jackson 2020 (TF-deletion single-cell atlas)}\newline Jackson CA, Castro DM, Saldi GA, Bonneau R, Gresham D. eLife 2020;9:e51254.\newline \href{https://doi.org/10.7554/eLife.51254}{doi:10.7554/\allowbreak eLife.51254} & Transcribed genotype barcodes with a per-cell transcriptome readout, fully crossed over 11 conditions. The design a yeast Perturb-seq is specified against. & GEO GSE125162; also eLife Source code 2 (103118\_SS\_Data.tsv.gz). Confirmed from the mirrored paper. The released matrix carries 38,225 cells; the abstract says 38,285 and the discussion 38,255, so the deposit governs. The authoritative 72-strain genotype list is Supplementary file 1 Table S2, an Excel file not held in the mirror. \\
\addlinespace[5pt]
7 & \textbf{Puddu 2019 (WGS of the deletion collection)}\newline Puddu F, et al., Jackson SP. Nature 2019;573:416-420.\newline \href{https://doi.org/10.1038/s41586-019-1549-9}{doi:10.1038/\allowbreak s41586-019-1549-9} & Whole-genome sequence for every strain in the deletion collection. A pooled campaign reads a barcode and infers a genotype; this is the measurement of how often that inference is wrong, and Jackson 2020 built its own strains rather than use the collection for related reasons. & ENA/\allowbreak SRA (accession unconfirmed) \\
\addlinespace[5pt]
8 & \textbf{N'Guessan 2025 (segregant scRNA-seq eQTL)}\newline N'Guessan A, Boocock J, Kruglyak L, Albert FW. eLife 2025. PMC12303567.\newline \href{https://pmc.ncbi.nlm.nih.gov/articles/PMC12303567/}{PMC12303567} & About 4,500 sequenced segregants profiled by single-cell RNA-seq, so genotype and transcriptome are measured in the same cell. & eLife data availability; segregant panel from Bloom lineage \\
\addlinespace[5pt]
9 & \textbf{Hackett 2020 (IDEA inducible-TF transcriptome time series)}\newline Hackett SR, Baltz EA, Coram M, Wranik BJ, Kim G, Baker A, Fan M, Hendrickson DG, Berndl M, McIsaac RS. Mol Syst Biol 2020;16:e9174.\newline \href{https://doi.org/10.15252/msb.20199174}{doi:10.15252/\allowbreak msb.20199174} & Induction rather than deletion, with a transcriptome followed over time. Jackson 2020 names transient induction as the perturbation modality most likely to produce a detectable expression response. & Published with the paper (Mol Syst Biol SI); accession unconfirmed from the mirror \\
\addlinespace[5pt]
10 & \textbf{Hu 2007 (TF deletion expression compendium)}\newline Hu Z, Killion PJ, Iyer VR. Nat Genet 2007;39:683-687.\newline \href{https://doi.org/10.1038/ng2012}{doi:10.1038/\allowbreak ng2012} & A designed single-gene perturbation crossed with a transcriptome over 269 transcription-factor deletions, which is the bulk form of what the campaign measures per cell, and the widest regulator coverage available. & GEO (unconfirmed) \\
\addlinespace[5pt]
11 & \textbf{Dong 2021 (MAGIC + SAM biosensor)}\newline Dong C, Schultz JC, Liu W, Lian J, Huang L, Xu Z, Zhao H. Metab Eng 2021;66:319-327.\newline \href{https://doi.org/10.1016/j.ymben.2021.03.005}{doi:10.1016/\allowbreak j.ymben.2021.03.005} & A genome-wide CRISPR activation, interference and deletion library sorted on a metabolite biosensor, so the label is product concentration rather than fitness. The pattern a sorted Perturb-seq would reuse. & Metab Eng SI (accession unconfirmed) \\
\addlinespace[5pt]
12 & \textbf{Momen-Roknabadi 2020 (inducible CRISPRi library)}\newline Momen-Roknabadi A, Oikonomou P, Zegans M, Tavazoie S. Commun Biol 2020;3:723.\newline \href{https://doi.org/10.1038/s42003-020-01452-9}{doi:10.1038/\allowbreak s42003-020-01452-9} & A genome-scale inducible CRISPR interference library with a released per-guide matrix. Induction control is what lets a knockdown be applied after a cell is captured rather than during outgrowth. & Addgene + Commun Biol SI \\
\addlinespace[5pt]
13 & \textbf{McGlincy 2021 (genome-scale CRISPRi library)}\newline McGlincy NJ, Meacham ZA, Reynaud KK, Muller R, Baum R, Ingolia NT. BMC Genomics 2021;22:205.\newline \href{https://doi.org/10.1186/s12864-021-07518-0}{doi:10.1186/\allowbreak s12864-021-07518-0} & The second independently designed genome-scale yeast CRISPR interference library with a released per-guide matrix, so guide design and gene effect can be separated. & ingolia-lab.org/\allowbreak yeast-crispri + Addgene + BMC SI \\
\addlinespace[5pt]
14 & \textbf{Bao 2018 (CHAnGE single-nucleotide library)}\newline Bao Z, HamediRad M, Xue P, Xiao H, Tasan I, Chao R, Liang J, Zhao H. Nat Biotechnol 2018;36:505-508.\newline \href{https://doi.org/10.1038/nbt.4132}{doi:10.1038/\allowbreak nbt.4132} & A guide-indexed genome-wide library whose edits sit below the open reading frame, so the perturbation is an allele rather than a null. & Nat Biotechnol SI (raw sequencing location unconfirmed) \\
\addlinespace[5pt]
15 & \textbf{Crook 2016 (tunable RNAi, isobutanol + 1-butanol)}\newline Crook N, Sun J, Morse N, Schmitz A, Alper HS. Appl Microbiol Biotechnol 2016;100:10005-10018. PMID 27654654.\newline \href{https://pubmed.ncbi.nlm.nih.gov/27654654/}{PMID 27654654} & A dose-graded knockdown axis. Graded rather than binary perturbation is what turns a per-cell readout into a dose-response curve. & AMB SI (DOI is an open verification item) \\
\addlinespace[5pt]
16 & \textbf{Roy 2018 (multiplexed precision editing)}\newline Roy KR, Smith JD, Vonesch SC, et al., St Onge RP. Nat Biotechnol 2018;36:512-520.\newline \href{https://doi.org/10.1038/nbt.4137}{doi:10.1038/\allowbreak nbt.4137} & Multiplexed designed edits per cell, which is combinatorial input in the sense the Perturb-seq specification requires. & Nat Biotechnol SI (accession unconfirmed) \\
\addlinespace[5pt]
17 & \textbf{Newman 2006 (protein noise, GFP library)}\newline Newman JRS, Ghaemmaghami S, Ihmels J, Breslow DK, Noble M, DeRisi JL, Weissman JS. Nature 2006;441:840-846.\newline \href{https://doi.org/10.1038/nature04785}{doi:10.1038/\allowbreak nature04785} & Per-cell protein abundance distributions across a genome-scale tagged collection. Expression noise per gene is what sets how large a perturbation effect has to be before a per-cell readout can resolve it. & Nature SI \\
\addlinespace[5pt]
18 & \textbf{Mukherjee 2021 (CRISPRi essential genes x acetic acid)}\newline Mukherjee V, Lind U, St Onge RP, Blomberg A, Nystrom T. mSystems 2021;6:e00410-21.\newline \href{https://doi.org/10.1128/mSystems.00410-21}{doi:10.1128/\allowbreak mSystems.00410-21} & Knockdown reaches the essential genes a deletion collection cannot hold, which is a third of the genome a deletion-based campaign would miss. & mSystems open-access SI \\
\addlinespace[5pt]
19 & \textbf{Lian 2017 (CRISPR-AID, beta-carotene)}\newline Lian J, HamediRad M, Hu S, Zhao H. Nat Commun 2017;8:1688.\newline \href{https://doi.org/10.1038/s41467-017-01695-x}{doi:10.1038/\allowbreak s41467-017-01695-x} & Three perturbation modalities in one cell, activation, interference and deletion, which is combinatorial input by modality rather than by gene. & Nat Commun SI \\
\addlinespace[5pt]
20 & \textbf{Hughes 2000 (compendium of expression profiles)}\newline Hughes TR, Marton MJ, Jones AR, Roberts CJ, Stoughton R, et al., Friend SH. Cell 2000;102:109-126.\newline \href{https://doi.org/10.1016/S0092-8674(00)00015-5}{doi:10.1016/\allowbreak S0092-8674(00)00015-5} & The original designed-perturbation expression compendium, and the independent measurement that makes cross-laboratory replication of the deletion transcriptome testable against the built Kemmeren 2014. & Rosetta compendium, distributed with the paper; not fetched \\
\addlinespace[5pt]
21 & \textbf{Lenstra 2011 (chromatin regulator deletion expression)}\newline Lenstra TL, Benschop JJ, Kim T, et al., Holstege FCP. Mol Cell 2011;42:536-549.\newline \href{https://doi.org/10.1016/j.molcel.2011.03.026}{doi:10.1016/\allowbreak j.molcel.2011.03.026} & The same design over chromatin regulators rather than sequence-specific factors, so the two together cover both halves of transcriptional control by deletion. & GEO (unconfirmed) \\
\addlinespace[5pt]
22 & \textbf{Sun 2013 (mRNA synthesis and decay rates across deletion strains)}\newline Sun M, Schwalb B, Pirkl N, Maier KC, Schenk A, Failmezger H, Tresch A, Cramer P. Mol Cell 2013;52:52-62.\newline \href{https://doi.org/10.1016/j.molcel.2013.09.010}{doi:10.1016/\allowbreak j.molcel.2013.09.010} & A designed single-gene perturbation crossed with a transcriptome, and the only one that resolves the level into a synthesis rate and a decay rate. Rates are what a per-cell snapshot cannot recover on its own. & GEO, series not fetched this pass \\
\addlinespace[5pt]
23 & \textbf{Jariani 2020 (yeast scRNA-seq through lag phase)}\newline Jariani A, Vermeersch L, Cerulus B, Perez-Samper G, Voordeckers K, Van Brussel T, Thienpont B, Lambrechts D, Verstrepen KJ. eLife 2020;9:e55320.\newline \href{https://doi.org/10.7554/eLife.55320}{doi:10.7554/\allowbreak eLife.55320} & Per-cell states through a carbon-source shift, so the readout is the distribution across an unsynchronized population rather than its mean. & GEO GSE144820 and GSE116246; from the collection sweep, not fetched \\
\addlinespace[5pt]
24 & \textbf{Guo 2018 (CRISPR-Cas9 tiling of essential genes)}\newline Guo X, Chavez A, Tung A, et al., Church GM, Shalem O. Nat Biotechnol 2018;36:540-546.\newline \href{https://doi.org/10.1038/nbt.4147}{doi:10.1038/\allowbreak nbt.4147} & Guide tiling produces a graded allelic series across one gene, so guide position rather than gene identity carries the effect. & Nat Biotechnol SI (unconfirmed) \\
\addlinespace[5pt]
25 & \textbf{Urbonaite 2021 (yeastDrop-Seq under drug treatment)}\newline Urbonaite G, Lee JTH, Liu P, Parras GG, Hemberg M, Acar M. Commun Biol 2021;4:1372.\newline \href{https://doi.org/10.1038/s42003-021-02895-4}{doi:10.1038/\allowbreak s42003-021-02895-4} & Single-drug arms against their combination with a per-cell readout, the environmental analog of a combinatorial genetic perturbation. & GEO GSE165686; Zenodo 10.5281/\allowbreak zenodo.4767298 and 10.5281/\allowbreak zenodo.4762526; from the collection sweep, not fetched \\
\addlinespace[5pt]
26 & \textbf{Nadal-Ribelles 2019 (sensitive yeast scRNA-seq)}\newline Nadal-Ribelles M, Islam S, Wei W, et al., Posas F, Steinmetz LM. Nat Microbiol 2019;4:683-692.\newline \href{https://doi.org/10.1038/s41564-018-0346-9}{doi:10.1038/\allowbreak s41564-018-0346-9} & The high-sensitivity yeast protocol, strand and isoform aware. Sensitivity per cell is the constraint on how small a perturbation effect a campaign can resolve. & Nat Microbiol SI /\allowbreak  GEO \\
\addlinespace[5pt]
27 & \textbf{Gasch 2017 (single-cell stress heterogeneity)}\newline Gasch AP, Yu FB, Hose J, et al., Quake SR. PLoS Biol 2017;15:e2004050.\newline \href{https://doi.org/10.1371/journal.pbio.2004050}{doi:10.1371/\allowbreak journal.pbio.2004050} & Separates intrinsic from extrinsic per-cell variation in an isogenic population under stress, which is the null a perturbation effect is measured against. & PLoS Biology SI /\allowbreak  GEO \\
\addlinespace[5pt]
28 & \textbf{Jaffe 2019 (multiplexed CRISPR interference epistasis)}\newline Jaffe M, Sherlock G, Levy SF. G3 2019;9:2843-2852.\newline \href{https://doi.org/10.1534/g3.119.400323}{doi:10.1534/\allowbreak g3.119.400323} & Two knockdowns in the same cell with a measured interaction, which is the only combinatorial CRISPR interference design located in yeast. & G3 open-access SI (unconfirmed) \\
\addlinespace[5pt]
29 & \textbf{Jackson 2023 (wild-type scRNA-seq time course for RNA kinetics)}\newline Jackson CA, Beheler-Amass M, Tjarnberg A, Suresh I, Hickey AS, Bonneau R, Gresham D. bioRxiv 2023. 10.1101/2023.09.21.558277.\newline \href{https://doi.org/10.1101/2023.09.21.558277}{doi:10.1101/\allowbreak 2023.09.21.558277} & The platform companion of Jackson 2020: 173,361 cells of one strain, which is what sets the per-cell variance a perturbation has to exceed. & GEO GSE242556 (from the scYeast Methods, Sec. 4.1.3) \\
\addlinespace[5pt]
30 & \textbf{Brettner 2024 (ultra-high-throughput yeast scRNA-seq)}\newline Brettner L, et al. Yeast 2024. doi:10.1002/yea.3927.\newline \href{https://doi.org/10.1002/yea.3927}{doi:10.1002/\allowbreak yea.3927} & Combinatorial barcoding in yeast at 96 to 384 multiplexed genotypes or environments per run. The throughput route by which a campaign becomes affordable. & Yeast (Wiley) SI \\
\addlinespace[5pt]
31 & \textbf{Mulleder 2012 (prototrophic deletion collection)}\newline Mulleder M, Capuano F, Pir P, et al., Ralser M. Nat Biotechnol 2012;30:1176-1178.\newline \href{https://doi.org/10.1038/nbt.2442}{doi:10.1038/\allowbreak nbt.2442} & The prototrophic deletion collection. Jackson 2020 used a prototrophic background because auxotrophy blocks minimal and nitrogen-limited media, so this is the strain resource a campaign in defined media needs. & EUROSCARF; Addgene kit 40276 \\
\addlinespace[5pt]
32 & \textbf{Su 2023 (single-cell transcriptomes under four stresses)}\newline Su Y, Xu C, Shea J, DeStephanis D, Su Z. BMC Genomics 2023;24:88.\newline \href{https://doi.org/10.1186/s12864-023-09184-w}{doi:10.1186/\allowbreak s12864-023-09184-w} & Full-length rather than three-prime tag counting, so isoform-level readout is on the table for a campaign that needs it. & GEO GSE201386 (from the scYeast Methods) \\
\addlinespace[5pt]
33 & \textbf{Wang 2022 (single-cell transcriptomes across replicative aging)}\newline Wang J, Sang Y, Jin S, Wang X, Azad GK, McCormick MA, Kennedy BK, Li Q, Wang J, Zhang X, et al. Aging Cell 2022;21:e13712.\newline \href{https://doi.org/10.1111/acel.13712}{doi:10.1111/\allowbreak acel.13712} & The only yeast single-cell set with age as the condition axis, which is a covariate any pooled campaign carries whether or not it measures it. & GEO GSE210032 (from the scYeast Methods) \\
\addlinespace[5pt]
34 & \textbf{Albert 2018 (eQTL in 1,012 segregants)}\newline Albert FW, Bloom JS, Siegel J, Day L, Kruglyak L. eLife 2018;7:e35471.\newline \href{https://doi.org/10.7554/eLife.35471}{doi:10.7554/\allowbreak eLife.35471} & The transcriptome half of the segregant panel that Jakobson 2025, Gerke 2017 and Eder 2020 measure protein, metabolite and flux on. & GEO /\allowbreak  eLife data availability (unconfirmed) \\
\addlinespace[5pt]
35 & \textbf{Jakobson 2025 (genome-to-proteome map)}\newline Jakobson CM, et al. Science 2025;390:eadu3198.\newline \href{https://doi.org/10.1126/science.adu3198}{doi:10.1126/\allowbreak science.adu3198} & Protein abundance across the same sequenced segregant panel Albert 2018 profiles by transcriptome, so protein and transcript quantitative trait loci are measurable on one genotype set. & Science data-availability statement (accession unconfirmed) \\
\addlinespace[5pt]
36 & \textbf{Aulakh 2025 (genome-scale ionome)}\newline Aulakh SK, et al. Cell Syst 2025;16:101319.\newline \href{https://doi.org/10.1016/j.cels.2025.101319}{doi:10.1016/\allowbreak j.cels.2025.101319} & A metabolic readout on the whole deletion collection, which is the genotype axis the supported expression compendium already covers. & Cell Systems SI (accession unconfirmed) \\
\addlinespace[5pt]
37 & \textbf{Teyssonniere 2024 (species-wide proteome against transcriptome)}\newline Teyssonniere EM, Trebulle P, Muenzner J, Loegler V, Ludwig D, Amari F, Mulleder M, Friedrich A, Hou J, Ralser M, Schacherer J. Proc Natl Acad Sci USA 2024;121:e2319211121.\newline \href{https://doi.org/10.1073/pnas.2319211121}{doi:10.1073/\allowbreak pnas.2319211121} & Proteomes for 942 isolates, the panel whose transcriptomes the built Caudal 2024 supplies, so the strain-level pairing is one to one. & PNAS SI; PMC11087752. PRIDE identifier not fetched this pass \\
\addlinespace[5pt]
38 & \textbf{Grossbach 2022 (BY x RM transcriptome, proteome and phosphoproteome)}\newline Grossbach J, Gillet L, Clement-Ziza M, Schmalohr CL, Schubert OT, Schutter M, Mawer JSP, Barnes CA, Bludau I, Weith M, Tessarz P, Graef M, Aebersold R, Beyer A. Mol Syst Biol 2022;18:e10712.\newline \href{https://doi.org/10.15252/msb.202110712}{doi:10.15252/\allowbreak msb.202110712} & Three layers on one set of 112 sequenced strains, so the transcript and the protein of a gene are observed in the same genotype and the RNA to protein question becomes a within-strain residual. & Mol Syst Biol SI + PRIDE, identifier not fetched this pass \\
\addlinespace[5pt]
39 & \textbf{Leutert 2023 (phosphoproteome x 101 conditions)}\newline Leutert M, et al., Villen J. Nat Struct Mol Biol 2023;30:1761-1773.\newline \href{https://doi.org/10.1038/s41594-023-01115-3}{doi:10.1038/\allowbreak s41594-023-01115-3} & Enzyme regulation acts faster than transcription, so a phosphosite layer over 101 conditions is what explains flux changes an expression table cannot. & PRIDE PXD034997 \\
\addlinespace[5pt]
40 & \textbf{Brauer 2008 (growth-rate-controlled chemostat transcriptome)}\newline Brauer MJ, Huttenhower C, Airoldi EM, et al., Botstein D. Mol Biol Cell 2008;19:352-367.\newline \href{https://doi.org/10.1091/mbc.e07-08-0779}{doi:10.1091/\allowbreak mbc.e07-08-0779} & Expression at controlled growth rate under each nutrient limitation, which is the condition axis Boer 2010 and Hackett 2016 measure metabolites and flux on. & GEO (unconfirmed) \\
\addlinespace[5pt]
41 & \textbf{Hackett 2016 (SIMMER multi-omic flux)}\newline Hackett SR, Zanotelli VRT, Xu W, et al., Rabinowitz JD. Science 2016;354:aaf2786.\newline \href{https://doi.org/10.1126/science.aaf2786}{doi:10.1126/\allowbreak science.aaf2786} & Flux, metabolite and transcript measured in one study under the same nutrient limitations, so the join is internal rather than across papers. & Science SI (accession unconfirmed) \\
\addlinespace[5pt]
42 & \textbf{Zhu 2014 (kinase / phosphatase lipidomics)}\newline Zhu Z, Loewen CJR. Mol Biol Cell 2014. PMC4196872.\newline \href{https://pmc.ncbi.nlm.nih.gov/articles/PMC4196872/}{PMC4196872} & A lipidome over 129 signaling deletions, most of which carry an expression profile in the supported compendium. & Mol Biol Cell SI Table S4 (Excel) \\
\addlinespace[5pt]
43 & \textbf{Gerke 2017 (urea-cycle mQTL)}\newline Gerke J, et al., Fay JC. Genetics 2017;206:2199.\newline \href{https://doi.org/10.1534/genetics.117.201107}{doi:10.1534/\allowbreak genetics.117.201107} & Metabolite quantitative trait loci on a segregant cross that also has an expression map, so a metabolite locus can be read through its transcript. & Genetics (GSA) SI \\
\addlinespace[5pt]
44 & \textbf{Ambroset 2014 (metabolite QTL)}\newline Ambroset C, et al., Fay JC. PLoS Genet 2014;10:e1004142.\newline \href{https://doi.org/10.1371/journal.pgen.1004142}{doi:10.1371/\allowbreak journal.pgen.1004142} & A 74-metabolite panel on a sequenced segregant panel, the widest metabolite vector available on a recombinant genotype axis. & PLOS Genetics SI \\
\addlinespace[5pt]
45 & \textbf{Blank 2005 (13C metabolic flux)}\newline Blank LM, Kuepfer L, Sauer U. Genome Biol 2005;6:R49.\newline \href{https://doi.org/10.1186/gb-2005-6-6-r49}{doi:10.1186/\allowbreak gb-2005-6-6-r49} & The only row measuring flux rather than a concentration, on deletion mutants that also have a transcriptome in the supported set. & Genome Biology open-access SI \\
\addlinespace[5pt]
46 & \textbf{Skelly 2013 (expression variation across isolates)}\newline Skelly DA, Merrihew GE, Riffle M, et al., MacCoss MJ, Akey JM. Genome Res 2013;23:1496-1504.\newline \href{https://doi.org/10.1101/gr.155762.113}{doi:10.1101/\allowbreak gr.155762.113} & Transcriptome and proteome on the same 22 strains, which is the only row where the cheap and expensive layers are paired within one experiment rather than joined across two. & Genome Research SI (unconfirmed) \\
\addlinespace[5pt]
47 & \textbf{Airoldi 2016 (nitrogen-limited steady-state and dynamic transcriptome)}\newline Airoldi EM, Miller D, Athanasiadou R, Brandt N, Abdul-Rahman F, Neymotin B, Hashimoto T, Bahmani T, Gresham D. Mol Biol Cell 2016;27:1383-1396.\newline \href{https://doi.org/10.1091/mbc.E14-05-1013}{doi:10.1091/\allowbreak mbc.E14-05-1013} & Expression under the exact nitrogen-limited conditions that Jackson 2020 profiles single cells in, so the deletion effect separates from the medium effect. & GEO, series not confirmed this pass \\
\addlinespace[5pt]
48 & \textbf{McManus 2014 (ribosome profiling, allele-specific)}\newline McManus CJ, May GE, Spealman P, Shteyman A. Genome Res 2014;24:422-430.\newline \href{https://doi.org/10.1101/gr.164996.113}{doi:10.1101/\allowbreak gr.164996.113} & Translation efficiency per gene alongside the mRNA abundance of that gene, which is the first of the two terms that separate protein level from transcript level. & GEO, series not fetched this pass \\
\addlinespace[5pt]
49 & \textbf{Martin-Perez 2017 (protein half-lives)}\newline Martin-Perez M, Villen J. Cell Syst 2017;5:283-294.e5.\newline \href{https://doi.org/10.1016/j.cels.2017.08.008}{doi:10.1016/\allowbreak j.cels.2017.08.008} & Protein turnover per gene, the second of those two terms, spanning three orders of magnitude and therefore not replaceable by a global constant. & Cell Systems SI; PRIDE identifier not fetched this pass \\
\addlinespace[5pt]
50 & \textbf{Boer 2010 (metabolome across nutrient limitations)}\newline Boer VM, Crutchfield CA, Bradley PH, Botstein D, Rabinowitz JD. Mol Biol Cell 2010;21:198-211.\newline \href{https://doi.org/10.1091/mbc.e09-07-0597}{doi:10.1091/\allowbreak mbc.e09-07-0597} & The metabolite half of the chemostat nutrient-limitation series whose transcriptome half is Brauer 2008, same laboratory and same conditions. & MBoC SI (unconfirmed) \\
\addlinespace[5pt]
51 & \textbf{Tengolics 2024 (domestication metabolome)}\newline Tengolics R, Szappanos B, Mulleder M, et al., Ralser M, Papp B. PNAS 2024;121:e2313354121.\newline \href{https://doi.org/10.1073/pnas.2313354121}{doi:10.1073/\allowbreak pnas.2313354121} & Metabolite levels across sequenced isolates, the panel the supported Caudal 2024 transcriptomes and Muenzner 2024 proteomes also sit on. & PNAS Dataset S11 \\
\addlinespace[5pt]
52 & \textbf{Eder 2020 (flux QTL)}\newline Eder M, Nidelet T, Sanchez I, Camarasa C, Legras JL, Dequin S. PMID 32034164.\newline \href{https://pubmed.ncbi.nlm.nih.gov/32034164/}{PMID 32034164} & Flux mapped to segregant genotypes, which is the flux counterpart of the expression and protein maps on the same panel type. & Journal SI (venue and accession unconfirmed) \\
\addlinespace[5pt]
53 & \textbf{Foss 2007 (BY x RM segregant proteome)}\newline Foss EJ, Radulovic D, Shaffer SA, Ruderfer DM, Bedalov A, Goodlett DR, Kruglyak L. Nat Genet 2007;39:1369-1375.\newline \href{https://doi.org/10.1038/ng.2007.22}{doi:10.1038/\allowbreak ng.2007.22} & The independent replicate of Grossbach 2022 on the same cross, and the first measurement that protein-level loci differ from transcript-level loci. & Nat Genet supplementary tables; not fetched \\
\addlinespace[5pt]
54 & \textbf{Yu 2021 (proteome and metabolome under nitrogen limitation)}\newline Yu R, Vorontsov E, Sihlbom C, Nielsen J. eLife 2021;10:e65722.\newline \href{https://doi.org/10.7554/eLife.65722}{doi:10.7554/\allowbreak eLife.65722} & Both layers measured in one study under nitrogen limitation, so it calibrates the protein-to-metabolite step that cross-study joins assume. & eLife data availability (unconfirmed) \\
\addlinespace[5pt]
55 & \textbf{de Boer 2020 (100M random promoters)}\newline de Boer CG, Vaishnav ED, Sadeh R, Abeyta EL, Friedman N, Regev A. Nat Biotechnol 2020;38:56-65.\newline \href{https://doi.org/10.1038/s41587-019-0315-8}{doi:10.1038/\allowbreak s41587-019-0315-8} & Scale band: tier 1, ordered by measurements (100,000,000). & GEO (de Boer 2020); DREAM 2022 release of 6,739,258 promoters \\
\addlinespace[5pt]
56 & \textbf{Vaishnav 2022 (regulatory DNA fitness landscape)}\newline Vaishnav ED, de Boer CG, Molinet J, Yassour M, Fan L, Adiconis X, Thompson DA, Levin JZ, Cubillos FA, Regev A. Nature 2022;603:455-463.\newline \href{https://doi.org/10.1038/s41586-022-04506-6}{doi:10.1038/\allowbreak s41586-022-04506-6} & Scale band: tier 1, ordered by measurements (20,000,000). & GEO per the Nature data-availability statement (unconfirmed) \\
\addlinespace[5pt]
57 & \textbf{Lee 2014 (HIP-HOP fitness signatures)}\newline Lee AY, St Onge RP, Proctor MJ, et al., Giaever G, Nislow C. Science 2014;344:208-211.\newline \href{https://doi.org/10.1126/science.1250217}{doi:10.1126/\allowbreak science.1250217} & Scale band: tier 1, ordered by measurements (13,000,000). & Science SI + Nislow/\allowbreak Giaever portal, neither of which yielded per-strain matrices; awaiting author matrices per the WS15 roadmap \\
\addlinespace[5pt]
58 & \textbf{Nguyen Ba 2022 (barcoded bulk QTL, 100k segregants)}\newline Nguyen Ba AN, Lawrence KR, Rego-Costa A, Gopalakrishnan S, Temko D, Michor F, Desai MM. eLife 2022;11:e73983.\newline \href{https://doi.org/10.7554/eLife.73983}{doi:10.7554/\allowbreak eLife.73983} & Scale band: tier 1, ordered by measurements (800,000). & eLife data availability; SRA (unconfirmed) \\
\addlinespace[5pt]
59 & \textbf{Galardini 2019 (four backgrounds x 38 conditions)}\newline Galardini M, Busby BP, Vieitez C, Dunham AS, Typas A, Beltrao P. Mol Syst Biol 2019;15:e8831.\newline \href{https://doi.org/10.15252/msb.20198831}{doi:10.15252/\allowbreak msb.20198831} & Scale band: tier 1, ordered by measurements (575,472). & GEO GSE123118 + github.com/\allowbreak mgalardini/\allowbreak 2018koyeast \\
\addlinespace[5pt]
60 & \textbf{Parsons 2006 (bioactive-compound profiling)}\newline Parsons AB, Lopez A, Givoni IE, et al., Boone C. Cell 2006;126:611-625.\newline \href{https://doi.org/10.1016/j.cell.2006.06.040}{doi:10.1016/\allowbreak j.cell.2006.06.040} & Scale band: tier 1, ordered by measurements (410,000). & Cell SI tables (accession unconfirmed) \\
\addlinespace[5pt]
61 & \textbf{Dutta 2026 (barcoded natural-isolate chemical response)}\newline Dutta A, Garin M, Loegler V, Brach G, Friedrich A, Yoshimura M, Hirano H, Osada H, Boone C, Yashiroda Y, Hou J, Schacherer J. Nat Commun 2026.\newline \href{https://doi.org/10.1038/s41467-026-73532-z}{doi:10.1038/\allowbreak s41467-026-73532-z} & Scale band: tier 1, ordered by measurements (312,000). & Nat Commun SI; isolates from the 1,011 panel (ENA) \\
\addlinespace[5pt]
62 & \textbf{Li 2016 (tRNA fitness landscape)}\newline Li C, Qian W, Maclean CJ, Zhang J. Science 2016;352:837-840.\newline \href{https://doi.org/10.1126/science.aae0568}{doi:10.1126/\allowbreak science.aae0568} & Scale band: tier 1, ordered by measurements (65,000). & Science SI (unconfirmed) \\
\addlinespace[5pt]
63 & \textbf{Domingo 2018 (tRNA double-mutant landscape)}\newline Domingo J, Diss G, Lehner B. Nature 2018;558:117-121.\newline \href{https://doi.org/10.1038/s41586-018-0170-7}{doi:10.1038/\allowbreak s41586-018-0170-7} & Scale band: tier 1, ordered by measurements (23,000). & Nature SI (unconfirmed) \\
\addlinespace[5pt]
64 & \textbf{Zhang 2021 (YETI titratable overexpression)}\newline Zhang Y, Ho Yee Chow M, Ling Y, et al., Andrews BJ. Mol Syst Biol 2021;17:e10321.\newline \href{https://doi.org/10.15252/msb.202010321}{doi:10.15252/\allowbreak msb.202010321} & Scale band: tier 1, ordered by measurements (22,760). & Mol Syst Biol SI (repository accession unconfirmed) \\
\addlinespace[5pt]
65 & \textbf{Peeters 2021 (fermentation-trait QTL atlas)}\newline Peeters B, Reijbroek F, Verbaet J, et al., Jarosz DF, Verstrepen KJ. 2021. PMID 34727964.\newline \href{https://pubmed.ncbi.nlm.nih.gov/34727964/}{PMID 34727964} & Scale band: tier 1, ordered by measurements (20,250). & Journal SI; raw reads at SRA/\allowbreak ENA (accession unconfirmed) \\
\addlinespace[5pt]
66 & \textbf{Van Leeuwen 2024 (short-chain organic acids)}\newline Shared and specific genetic determinants of tolerance to acetic, butyric and octanoic acid. PMC10903034.\newline \href{https://www.ncbi.nlm.nih.gov/pmc/articles/PMC10903034/}{PMC10903034} & Scale band: tier 1, ordered by measurements (14,400). & Journal SI (open-access PMC) \\
\addlinespace[5pt]
67 & \textbf{Chica 2026 (AutoDRY autophagy screen)}\newline Chica N, et al. Nat Cell Biol 2026;28:465-479.\newline \href{https://doi.org/10.1038/s41556-025-01837-0}{doi:10.1038/\allowbreak s41556-025-01837-0} & Scale band: tier 1, ordered by measurements (11,838). & Dryad 10.5061/\allowbreak dryad.cfxpnvxdh \\
\addlinespace[5pt]
68 & \textbf{Sopko 2006 (genome-scale overexpression toxicity)}\newline Sopko R, Huang D, Preston N, et al., Boone C, Andrews B. Mol Cell 2006;21:319-330.\newline \href{https://doi.org/10.1016/j.molcel.2005.12.011}{doi:10.1016/\allowbreak j.molcel.2005.12.011} & Scale band: tier 1, ordered by measurements (10,560). & Mol Cell SI tables \\
\addlinespace[5pt]
69 & \textbf{Liu 2021 (tryptophan / isobutanol tolerance)}\newline Liu H-L, Wang CH-T, Chiang EP-I, Huang C-C, Li W-H. Biotechnol Biofuels 2021;14:200.\newline \href{https://doi.org/10.1186/s13068-021-02048-z}{doi:10.1186/\allowbreak s13068-021-02048-z} & Scale band: tier 1, ordered by measurements (10,012). & GEO GSE175794 (companion RNA-seq); screen in BMC SI \\
\addlinespace[5pt]
70 & \textbf{Kuroda 2019 (isobutanol-specific tolerance)}\newline Kuroda K, Hammer SK, Watanabe Y, Montano Lopez J, Fink GR, Stephanopoulos G, Ueda M, Avalos JL. Cell Syst 2019;9:534-547.e5.\newline \href{https://doi.org/10.1016/j.cels.2019.10.006}{doi:10.1016/\allowbreak j.cels.2019.10.006} & Scale band: tier 1, ordered by measurements (9,600). & ArrayExpress E-MTAB-8175 (RNA-seq); screen data in Cell Syst SI \\
\addlinespace[5pt]
71 & \textbf{Turco 2023 (Yeast Phenome)}\newline Turco G, Chang C, Wang RY, et al., Boone C, Andrews BJ, Roth FP. Sci Adv 2023;9:eadg5702.\newline \href{https://doi.org/10.1126/sciadv.adg5702}{doi:10.1126/\allowbreak sciadv.adg5702} & Scale band: tier 2, ordered by measurements (37,680,000). & yeastphenome.org (Zenodo 10.5281/\allowbreak zenodo.7714347); loader torchcell/\allowbreak datasets/\allowbreak scerevisiae/\allowbreak yeastphenome.py, 49 screens, not yet in the built set \\
\addlinespace[5pt]
72 & \textbf{Rossi 2021 (ChIP-exo protein architecture)}\newline Rossi MJ, Kuntala PK, Lai WKM, et al., Pugh BF. Nature 2021;592:309-314.\newline \href{https://doi.org/10.1038/s41586-021-03314-8}{doi:10.1038/\allowbreak s41586-021-03314-8} & Scale band: tier 2, ordered by measurements (2,400,000). & yeastepigenome.org; github.com/\allowbreak CEGRcode/\allowbreak 2021-Rossi\_Nature \\
\addlinespace[5pt]
73 & \textbf{Piotrowski 2017 (MOSAIC diagnostic panel)}\newline Piotrowski JS, Li SC, Deshpande R, et al., Boone C, Myers CL. Nat Chem Biol 2017;13:982-993.\newline \href{https://doi.org/10.1038/nchembio.2436}{doi:10.1038/\allowbreak nchembio.2436} & Scale band: tier 2, ordered by measurements (2,123,268). & MOSAIC portal (mosaic.cs.umn.edu) + Nat Chem Biol SI \\
\addlinespace[5pt]
74 & \textbf{Bloom 2025 (global epistasis in a yeast cross)}\newline Bloom JS, et al., Kruglyak L. 2025. PMID 40679398.\newline \href{https://pubmed.ncbi.nlm.nih.gov/40679398/}{PMID 40679398} & Scale band: tier 2, ordered by measurements (1,000,000). & journal SI (unconfirmed; recent) \\
\addlinespace[5pt]
75 & \textbf{Gasch 2000 (environmental stress response)}\newline Gasch AP, Spellman PT, Kao CM, et al., Brown PO. Mol Biol Cell 2000;11:4241-4257.\newline \href{https://doi.org/10.1091/mbc.11.12.4241}{doi:10.1091/\allowbreak mbc.11.12.4241} & Scale band: tier 2, ordered by measurements (900,000). & SPELL /\allowbreak  SGD Expression Connection; Stanford legacy mirror \\
\addlinespace[5pt]
76 & \textbf{Decourty 2021 (RNA-metabolism genetic interactions)}\newline Decourty L, et al., Saveanu C. Nucleic Acids Res 2021;49:8535-8555.\newline \href{https://doi.org/10.1093/nar/gkab680}{doi:10.1093/\allowbreak nar/\allowbreak gkab680} & Scale band: tier 2, ordered by measurements (700,000). & NAR open-access SI (accession unconfirmed) \\
\addlinespace[5pt]
77 & \textbf{Cuperus 2017 (random 5' UTR library)}\newline Cuperus JT, Groves B, Kuchina A, Rosenberg AB, Jojic N, Fields S, Seelig G. Genome Res 2017;27:2015-2024.\newline \href{https://doi.org/10.1101/gr.224964.117}{doi:10.1101/\allowbreak gr.224964.117} & Scale band: tier 2, ordered by measurements (500,000). & GEO (unconfirmed) \\
\addlinespace[5pt]
78 & \textbf{Faure 2022 (doubledeepPCA)}\newline Faure AJ, Domingo J, Schmiedel JM, Hidalgo-Carcedo C, Diss G, Lehner B. Nature 2022;604:175-183.\newline \href{https://doi.org/10.1038/s41586-022-04586-4}{doi:10.1038/\allowbreak s41586-022-04586-4} & Scale band: tier 2, ordered by measurements (500,000). & Nature data availability (unconfirmed) \\
\addlinespace[5pt]
79 & \textbf{Chong 2015 / CYCLoPs (single-cell proteome dynamics)}\newline Chong YT, Koh JLY, Friesen H, et al., Andrews BJ, Boone C. Cell 2015;161:1413-1424.\newline \href{https://doi.org/10.1016/j.cell.2015.04.051}{doi:10.1016/\allowbreak j.cell.2015.04.051} & Scale band: tier 2, ordered by measurements (418,200). & thecellvision.org/\allowbreak cyclops (bulk CSV) \\
\addlinespace[5pt]
80 & \textbf{Teyssonniere 2024 (species-wide trait survey)}\newline Teyssonniere EM, Shichino Y, Mito M, Friedrich A, Iwasaki S, Schacherer J. PLoS Genet 2024;20:e1011119.\newline \href{https://doi.org/10.1371/journal.pgen.1011119}{doi:10.1371/\allowbreak journal.pgen.1011119} & Scale band: tier 2, ordered by measurements (202,200). & PLoS Genetics SI (unconfirmed) \\
\addlinespace[5pt]
81 & \textbf{Diss 2018 (Fos-Jun interaction double mutants)}\newline Diss G, Lehner B. eLife 2018;7:e32472.\newline \href{https://doi.org/10.7554/eLife.32472}{doi:10.7554/\allowbreak eLife.32472} & Scale band: tier 2, ordered by measurements (120,000). & eLife data availability (unconfirmed) \\
\addlinespace[5pt]
82 & \textbf{Warringer 2011 (trait variation across the SGRP panel)}\newline Warringer J, Zorgo E, Cubillos FA, et al., Blomberg A, Liti G. PLoS Genet 2011;7:e1002111.\newline \href{https://doi.org/10.1371/journal.pgen.1002111}{doi:10.1371/\allowbreak journal.pgen.1002111} & Scale band: tier 2, ordered by measurements (120,000). & PLoS Genetics SI (unconfirmed) \\
\addlinespace[5pt]
83 & \textbf{Keren 2013 (promoter activity across conditions)}\newline Keren L, Zackay O, Lotan-Pompan M, Barenholz U, Dekel E, Sasson V, Aidelberg G, Bren A, Zeevi D, Weinberger A, Alon U, Milo R, Segal E. Mol Syst Biol 2013;9:701.\newline \href{https://doi.org/10.1038/msb.2013.59}{doi:10.1038/\allowbreak msb.2013.59} & Scale band: tier 2, ordered by measurements (90,000). & Mol Syst Biol SI (unconfirmed) \\
\addlinespace[5pt]
84 & \textbf{Puchta 2016 (snoRNA U3 landscape)}\newline Puchta O, Cseke B, Czaja H, Tollervey D, Sanguinetti G, Kudla G. Science 2016;352:840-844.\newline \href{https://doi.org/10.1126/science.aaf0965}{doi:10.1126/\allowbreak science.aaf0965} & Scale band: tier 2, ordered by measurements (60,000). & Science SI (unconfirmed) \\
\addlinespace[5pt]
85 & \textbf{Yoshikawa 2011 (deletion + overexpression phenome)}\newline Yoshikawa K, Tanaka T, Furusawa C, Nagahisa K, Hirasawa T, Shimizu H. G3 2011;1:247-267.\newline \href{https://doi.org/10.1534/g3.111.000695}{doi:10.1534/\allowbreak g3.111.000695} & Scale band: tier 2, ordered by measurements (57,600). & G3 open-access SI \\
\addlinespace[5pt]
86 & \textbf{Ho 2009 (MoBY-ORF barcoded overexpression)}\newline Ho CH, Magtanong L, Barker SL, et al., Boone C. Nat Biotechnol 2009;27:369-377.\newline \href{https://doi.org/10.1038/nbt.1534}{doi:10.1038/\allowbreak nbt.1534} & Scale band: tier 2, ordered by measurements (51,000). & moby.ccbr.utoronto.ca + Andrews lab portal \\
\addlinespace[5pt]
87 & \textbf{Bloom 2013 (1,008 BYxRM segregants)}\newline Bloom JS, Ehrenreich IM, Loo WT, Lite TL, Kruglyak L. Nature 2013;494:234-237.\newline \href{https://doi.org/10.1038/nature11867}{doi:10.1038/\allowbreak nature11867} & Scale band: tier 2, ordered by measurements (46,368). & Nature SI \\
\addlinespace[5pt]
88 & \textbf{Xiao 2014 (genome-wide RNAi furfural tolerance)}\newline Xiao H, Zhao H. Biotechnol Biofuels 2014;7:78.\newline \href{https://doi.org/10.1186/1754-6834-7-78}{doi:10.1186/\allowbreak 1754-6834-7-78} & Scale band: tier 2, ordered by measurements (20,000). & BMC open-access SI \\
\addlinespace[5pt]
89 & \textbf{Douglas 2012 (barcoded overexpression fitness)}\newline Douglas AC, Smith AM, Sharifpoor S, et al., Andrews BJ, Boone C, Nislow C. G3 2012;2:1279-1289.\newline \href{https://doi.org/10.1534/g3.112.003400}{doi:10.1534/\allowbreak g3.112.003400} & Scale band: tier 2, ordered by measurements (20,000). & G3 open-access SI \\
\addlinespace[5pt]
90 & \textbf{She 2018 (natural variant effects across a large isolate panel)}\newline She R, Jarosz DF. Cell 2018;172:478-490.\newline \href{https://doi.org/10.1016/j.cell.2017.12.015}{doi:10.1016/\allowbreak j.cell.2017.12.015} & Scale band: tier 2, ordered by measurements (20,000). & Cell SI (unconfirmed) \\
\addlinespace[5pt]
91 & \textbf{Costello 2020 (bioreactor Bar-seq)}\newline Costello Z, Wehrs M, Mukhopadhyay A, et al. Microb Cell Fact 2020;19:167.\newline \href{https://doi.org/10.1186/s12934-020-01423-z}{doi:10.1186/\allowbreak s12934-020-01423-z} & Scale band: tier 2, ordered by measurements (19,200). & Microb Cell Fact open-access SI \\
\addlinespace[5pt]
92 & \textbf{Gurvitz (fatty-acid utilization screen)}\newline Gurvitz A, et al. Mol Syst Biol (EMBO Press).\newline \href{https://www.embopress.org/doi/pdf/10.1038/msb4100051}{embopress.org/\allowbreak doi/\allowbreak pdf/\allowbreak 10.1038/\allowbreak msb4100051} & Scale band: tier 2, ordered by measurements (14,400). & Mol Syst Biol Table I (PDF-embedded) \\
\addlinespace[5pt]
93 & \textbf{Mira 2010 (acetic acid tolerance, full collection)}\newline Mira NP, Palma M, Guerreiro JF, Sa-Correia I. Microb Cell Fact 2010;9:79.\newline \href{https://doi.org/10.1186/1475-2859-9-79}{doi:10.1186/\allowbreak 1475-2859-9-79} & Scale band: tier 2, ordered by measurements (14,400). & Microb Cell Fact SI (PMC2972246) \\
\addlinespace[5pt]
94 & \textbf{Gameiro 2025 (AID2 degron libraries)}\newline Gameiro E, Juarez-Nunez KA, Fung JJ, Shankar S, Luke B, Khmelinskii A. J Cell Biol 2025;224:e202409007.\newline \href{https://doi.org/10.1083/jcb.202409007}{doi:10.1083/\allowbreak jcb.202409007} & Scale band: tier 2, ordered by measurements (11,200). & J Cell Biol SI; strains via IMB Mainz /\allowbreak  EUROSCARF \\
\addlinespace[5pt]
95 & \textbf{Si 2017 (RAGE RNAi, isobutanol titer)}\newline Si T, Chao R, Min Y, Wu Y, Ren W, Zhao H. Nat Commun 2017;8:15187.\newline \href{https://doi.org/10.1038/ncomms15187}{doi:10.1038/\allowbreak ncomms15187} & Scale band: tier 2, ordered by measurements (10,000). & Nat Commun SI (accession unconfirmed) \\
\addlinespace[5pt]
96 & \textbf{HamediRad 2018 (RAGE xylose utilization)}\newline HamediRad M, Lian J, Li H, Zhao H. Biotechnol Bioeng 2018;115:1552-1560.\newline \href{https://doi.org/10.1002/bit.26570}{doi:10.1002/\allowbreak bit.26570} & Scale band: tier 2, ordered by measurements (10,000). & Biotechnol Bioeng SI (accession unconfirmed) \\
\addlinespace[5pt]
97 & \textbf{Pereira 2014 (wheat-straw hydrolysate)}\newline Pereira FB, Guimaraes PMR, Gomes DG, et al., Domingues L. J Ind Microbiol Biotechnol 2014;41:1753-1761.\newline \href{https://doi.org/10.1007/s10295-014-1519-z}{doi:10.1007/\allowbreak s10295-014-1519-z} & Scale band: tier 2, ordered by measurements (9,600). & Journal SI \\
\addlinespace[5pt]
98 & \textbf{Endo 2008 (vanillin tolerance)}\newline Endo A, Nakamura T, Ando A, Tokuyasu K, Shima J. Appl Environ Microbiol 2008;74:7175-7185.\newline \href{https://doi.org/10.1128/AEM.01541-08}{doi:10.1128/\allowbreak AEM.01541-08} & Scale band: tier 2, ordered by measurements (9,600). & AEM SI (PMC2375868) \\
\addlinespace[5pt]
99 & \textbf{Sharifpoor 2012 (kinome synthetic dosage lethality)}\newline Sharifpoor S, van Dyk D, Costanzo M, et al., Boone C, Andrews BJ. Genome Res 2012;22:791-801.\newline \href{https://doi.org/10.1101/gr.129213.111}{doi:10.1101/\allowbreak gr.129213.111} & Scale band: tier 2, ordered by measurements (8,700). & Genome Research SI \\
\addlinespace[5pt]
100 & \textbf{Mavor 2016 (ubiquitin DMS across conditions)}\newline Mavor D, Barlow K, Thompson S, et al., Bolon DNA, Fraser JS. eLife 2016;5:e15802.\newline \href{https://doi.org/10.7554/eLife.15802}{doi:10.7554/\allowbreak eLife.15802} & Scale band: tier 2, ordered by measurements (7,000). & eLife data availability (unconfirmed) \\
\addlinespace[5pt]
101 & \textbf{Flynn 2020 (Hsp90 DMS across environments)}\newline Flynn JM, Rossouw A, Cote-Hammarlof P, et al., Bolon DNA. eLife 2020;9:e53810.\newline \href{https://doi.org/10.7554/eLife.53810}{doi:10.7554/\allowbreak eLife.53810} & Scale band: tier 2, ordered by measurements (6,000). & eLife data availability (unconfirmed) \\
\addlinespace[5pt]
102 & \textbf{Ho 2018 (unified absolute protein abundance)}\newline Ho B, Baryshnikova A, Brown GW. Cell Syst 2018;6:192-205.\newline \href{https://doi.org/10.1016/j.cels.2017.12.004}{doi:10.1016/\allowbreak j.cels.2017.12.004} & Scale band: tier 2, ordered by measurements (5,858). & Cell Systems SI \\
\addlinespace[5pt]
103 & \textbf{Renganaath 2020 (natural promoter-variant MPRA)}\newline Renganaath K, Cheung R, Day L, Kosuri S, Kruglyak L, Albert FW. eLife 2020;9:e62669.\newline \href{https://doi.org/10.7554/eLife.62669}{doi:10.7554/\allowbreak eLife.62669} & Scale band: tier 2, ordered by measurements (5,832). & eLife data availability; GEO (unconfirmed) \\
\addlinespace[5pt]
104 & \textbf{Chang 2012 (organic-acid production screen)}\newline Chang HJ, Suga H, et al. 2012. PMID 22277779.\newline \href{https://pubmed.ncbi.nlm.nih.gov/22277779/}{PMID 22277779} & Scale band: tier 2, ordered by measurements (4,800). & PDF tables only; manual transcription required \\
\addlinespace[5pt]
105 & \textbf{Trikka 2015 (carotenogenic heterozygous screen)}\newline Trikka FA, Nikolaidis A, Athanasakoglou A, et al., Kampranis SC, Makris AM. Microb Cell Fact 2015;14:60.\newline \href{https://doi.org/10.1186/s12934-015-0246-0}{doi:10.1186/\allowbreak s12934-015-0246-0} & Scale band: tier 2, ordered by measurements (4,700). & BMC Additional files 1-2 \\
\addlinespace[5pt]
106 & \textbf{Sardi 2018 (natural variation in hydrolysate tolerance)}\newline Sardi M, Paithane V, Place M, Robinson DE, Hose J, Wohlbach DJ, Gasch AP. PLoS Genet 2018;14:e1007217.\newline \href{https://doi.org/10.1371/journal.pgen.1007217}{doi:10.1371/\allowbreak journal.pgen.1007217} & Scale band: tier 2, ordered by measurements (2,000). & PLoS Genetics SI (unconfirmed) \\
\addlinespace[5pt]
107 & \textbf{Breslow 2008 (DAmP hypomorphic collection)}\newline Breslow DK, Cameron DM, Collins SR, et al., Weissman JS. Nat Methods 2008;5:711-718.\newline \href{https://doi.org/10.1038/nmeth.1234}{doi:10.1038/\allowbreak nmeth.1234} & Scale band: tier 2, ordered by measurements (1,812). & Nat Methods SI; Horizon YSC5050/\allowbreak 5090/\allowbreak 5093/\allowbreak 5094 \\
\addlinespace[5pt]
108 & \textbf{Kofoed 2015 (barcoded ts alleles)}\newline Kofoed M, Milbury KL, Chiang JH, et al., Hieter P, Stirling PC. G3 2015;5:1879-1887.\newline \href{https://doi.org/10.1534/g3.115.019174}{doi:10.1534/\allowbreak g3.115.019174} & Scale band: tier 2, ordered by measurements (1,200). & G3 open-access SI; strains via Dharmacon/\allowbreak Horizon \\
\addlinespace[5pt]
109 & \textbf{Zhang 2022 (GCE-SCRaMbLE recombination outcomes)}\newline Zhang W, Lazar-Stefanita L, Yamashita H, et al., Boeke JD. Nat Commun 2022;13:5836.\newline \href{https://doi.org/10.1038/s41467-022-33606-0}{doi:10.1038/\allowbreak s41467-022-33606-0} & Scale band: tier 2, ordered by measurements (200). & mirrored as zhangSystematicDissectionKey2022; SRA (unconfirmed) \\
\addlinespace[5pt]
110 & \textbf{Cubillos 2017 (nitrogen-consumption QTL)}\newline Cubillos FA, Brice C, Molinet J, et al., Martinez C. G3 2017;7:1693-1705.\newline \href{https://doi.org/10.1534/g3.117.042127}{doi:10.1534/\allowbreak g3.117.042127} & Scale band: tier 2, ordered by measurements (165). & BioProject PRJNA379146 (RNA-seq); genotypes in Cubillos 2013 Table S1 \\
\addlinespace[5pt]
111 & \textbf{Hillenmeyer 2010 (systematic fitness reanalysis)}\newline Hillenmeyer ME, Ericson E, Davis RW, Nislow C, Koller D, Giaever G. Genome Biol 2010;11:R30.\newline \href{https://doi.org/10.1186/gb-2010-11-3-r30}{doi:10.1186/\allowbreak gb-2010-11-3-r30} & Scale band: tier 3, ordered by measurements (2,360,000). & Genome Biology SI (unconfirmed) \\
\addlinespace[5pt]
112 & \textbf{Smith 2008 (gene-environment interaction eQTL)}\newline Smith EN, Kruglyak L. PLoS Biol 2008;6:e83.\newline \href{https://doi.org/10.1371/journal.pbio.0060083}{doi:10.1371/\allowbreak journal.pbio.0060083} & Scale band: tier 3, ordered by measurements (1,308,000). & GEO (unconfirmed) \\
\addlinespace[5pt]
113 & \textbf{Braberg 2020 (point-mutant E-MAP)}\newline Braberg H, et al., Krogan NJ. Science 2020;370:eaaz4910.\newline \href{https://doi.org/10.1126/science.aaz4910}{doi:10.1126/\allowbreak science.aaz4910} & Scale band: tier 3, ordered by measurements (500,000). & Science SI (accession unconfirmed) \\
\addlinespace[5pt]
114 & \textbf{Breker 2013 (GFP library under stress)}\newline Breker M, Gymrek M, Schuldiner M. J Cell Biol 2013;200:839-850.\newline \href{https://doi.org/10.1083/jcb.201301120}{doi:10.1083/\allowbreak jcb.201301120} & Scale band: tier 3, ordered by measurements (212,109). & LoQAtE (weizmann.ac.il/\allowbreak molgen/\allowbreak loqate) \\
\addlinespace[5pt]
115 & \textbf{Tkach 2012 (protein relocalization under DNA damage)}\newline Tkach JM, Yimit A, Lee AY, et al., Andrews BJ, Brown GW. Nat Cell Biol 2012;14:966-976.\newline \href{https://doi.org/10.1038/ncb2549}{doi:10.1038/\allowbreak ncb2549} & Scale band: tier 3, ordered by measurements (204,000). & TheCellVision /\allowbreak  Nat Cell Biol SI (unconfirmed) \\
\addlinespace[5pt]
116 & \textbf{Ehrenreich 2010 (X-QTL bulk segregant mapping)}\newline Ehrenreich IM, Torabi N, Jia Y, Kent J, Martis S, Shapiro JA, Gresham D, Caudy AA, Kruglyak L. Nature 2010;464:1039-1042.\newline \href{https://doi.org/10.1038/nature08923}{doi:10.1038/\allowbreak nature08923} & Scale band: tier 3, ordered by measurements (187,000). & Nature SI (unconfirmed) \\
\addlinespace[5pt]
117 & \textbf{Denervaud 2013 (microfluidic GFP dynamics)}\newline Denervaud N, Becker J, Delgado-Gonzalo R, et al., Maerkl SJ. PNAS 2013;110:15842-15847.\newline \href{https://doi.org/10.1073/pnas.1308265110}{doi:10.1073/\allowbreak pnas.1308265110} & Scale band: tier 3, ordered by measurements (68,000). & PNAS SI (unconfirmed) \\
\addlinespace[5pt]
118 & \textbf{Brogaard 2012 (chemical nucleosome map)}\newline Brogaard K, Xi L, Wang JP, Widom J. Nature 2012;486:496-501.\newline \href{https://doi.org/10.1038/nature11142}{doi:10.1038/\allowbreak nature11142} & Scale band: tier 3, ordered by measurements (67,000). & GEO (unconfirmed) \\
\addlinespace[5pt]
119 & \textbf{Xia 2022 (proteome allocation across conditions)}\newline Xia J, Sanchez BJ, Chen Y, Campbell K, Kasvandik S, Nielsen J. Nat Commun 2022;13:2819.\newline \href{https://doi.org/10.1038/s41467-022-30513-2}{doi:10.1038/\allowbreak s41467-022-30513-2} & Scale band: tier 3, ordered by measurements (60,000). & Nat Commun SI (unconfirmed) \\
\addlinespace[5pt]
120 & \textbf{Melamed 2013 (PAB1 RRM domain scan)}\newline Melamed D, Young DL, Gamble CE, Miller CR, Fields S. RNA 2013;19:1537-1551.\newline \href{https://doi.org/10.1261/rna.040709.113}{doi:10.1261/\allowbreak rna.040709.113} & Scale band: tier 3, ordered by measurements (40,000). & journal SI (unconfirmed) \\
\addlinespace[5pt]
121 & \textbf{Yvert 2013 (single-cell trait variation across isolates)}\newline Yvert G, Ohnuki S, Nogami S, Imanaga Y, Fehrmann S, Schacherer J, Ohya Y. BMC Syst Biol 2013;7:54.\newline \href{https://doi.org/10.1186/1752-0509-7-54}{doi:10.1186/\allowbreak 1752-0509-7-54} & Scale band: tier 3, ordered by measurements (18,537). & BMC SI (unconfirmed) \\
\addlinespace[5pt]
122 & \textbf{Lahtvee 2017 (absolute proteome across conditions)}\newline Lahtvee PJ, Sanchez BJ, Smialowska A, et al., Nielsen J. Cell Syst 2017;4:495-504.\newline \href{https://doi.org/10.1016/j.cels.2017.03.003}{doi:10.1016/\allowbreak j.cels.2017.03.003} & Scale band: tier 3, ordered by measurements (18,000). & Cell Systems SI (unconfirmed) \\
\addlinespace[5pt]
123 & \textbf{Weill 2018 (SWAT libraries)}\newline Weill U, Yofe I, Sass E, et al., Schuldiner M. Nat Methods 2018;15:617-622.\newline \href{https://doi.org/10.1038/s41592-018-0044-9}{doi:10.1038/\allowbreak s41592-018-0044-9} & Scale band: tier 3, ordered by measurements (16,500). & Nat Methods SI; strains via EUROSCARF \\
\addlinespace[5pt]
124 & \textbf{Kinsler 2020 (barcoded pleiotropy across environments)}\newline Kinsler G, Geiler-Samerotte K, Petrov DA. eLife 2020;9:e61271.\newline \href{https://doi.org/10.7554/eLife.61271}{doi:10.7554/\allowbreak eLife.61271} & Scale band: tier 3, ordered by measurements (13,500). & eLife data availability (unconfirmed) \\
\addlinespace[5pt]
125 & \textbf{Pelechano 2013 (transcript isoform landscape)}\newline Pelechano V, Wei W, Steinmetz LM. Nature 2013;497:127-131.\newline \href{https://doi.org/10.1038/nature12121}{doi:10.1038/\allowbreak nature12121} & Scale band: tier 3, ordered by measurements (12,000). & ArrayExpress /\allowbreak  GEO (unconfirmed) \\
\addlinespace[5pt]
126 & \textbf{Deutschbauer 2005 (haploinsufficiency, genome-wide)}\newline Deutschbauer AM, Jaramillo DF, Proctor M, Kumm J, Hillenmeyer ME, Davis RW, Nislow C, Giaever G. Genetics 2005;169:1915-1925.\newline \href{https://doi.org/10.1534/genetics.104.036871}{doi:10.1534/\allowbreak genetics.104.036871} & Scale band: tier 3, ordered by measurements (11,800). & Genetics SI (unconfirmed) \\
\addlinespace[5pt]
127 & \textbf{Valenti 2025 (SWAT degron + GFP collection)}\newline Valenti M, et al. J Cell Biol 2025;224:e202409050.\newline \href{https://doi.org/10.1083/jcb.202409050}{doi:10.1083/\allowbreak jcb.202409050} & Scale band: tier 3, ordered by measurements (10,000). & J Cell Biol SI \\
\addlinespace[5pt]
128 & \textbf{Sousa 2013 (acetate resistance overexpression screen)}\newline Sousa M, Duarte AM, Fernandes TR, et al. 2013. PMID 23262128.\newline \href{https://pubmed.ncbi.nlm.nih.gov/23262128/}{PMID 23262128} & Scale band: tier 3, ordered by measurements (10,000). & journal SI (unconfirmed) \\
\addlinespace[5pt]
129 & \textbf{Sadhu 2016 (CRISPR-directed causal variant mapping)}\newline Sadhu MJ, Bloom JS, Day L, Kruglyak L. Science 2016;352:1113-1116.\newline \href{https://doi.org/10.1126/science.aaf5124}{doi:10.1126/\allowbreak science.aaf5124} & Scale band: tier 3, ordered by measurements (10,000). & Science SI (unconfirmed) \\
\addlinespace[5pt]
130 & \textbf{Weinberg 2016 (ribosome profiling, improved)}\newline Weinberg DE, Shah P, Eichhorn SW, Hussmann JA, Plotkin JB, Bartel DP. Cell Rep 2016;14:1787-1799.\newline \href{https://doi.org/10.1016/j.celrep.2016.01.043}{doi:10.1016/\allowbreak j.celrep.2016.01.043} & Scale band: tier 3, ordered by measurements (6,000). & GEO (unconfirmed) \\
\addlinespace[5pt]
131 & \textbf{Ghaemmaghami 2003 (absolute protein abundance)}\newline Ghaemmaghami S, Huh WK, Bower K, et al., Weissman JS. Nature 2003;425:737-741.\newline \href{https://doi.org/10.1038/nature02046}{doi:10.1038/\allowbreak nature02046} & Scale band: tier 3, ordered by measurements (4,251). & Nature SI (Excel) \\
\addlinespace[5pt]
132 & \textbf{Huh 2003 (GFP localization atlas)}\newline Huh WK, Falvo JV, Gerke LC, et al., O'Shea EK. Nature 2003;425:686-691.\newline \href{https://doi.org/10.1038/nature02026}{doi:10.1038/\allowbreak nature02026} & Scale band: tier 3, ordered by measurements (4,156). & SGD /\allowbreak  YeastGFP mirror \\
\addlinespace[5pt]
133 & \textbf{Michaelis 2023 (protein interactome)}\newline Michaelis AC, et al., Mann M. Nature 2023;624:192-200.\newline \href{https://doi.org/10.1038/s41586-023-06739-5}{doi:10.1038/\allowbreak s41586-023-06739-5} & Scale band: tier 3, ordered by measurements (4,000). & yeast-interactome.biochem.mpg.de + ProteomeXchange \\
\addlinespace[5pt]
134 & \textbf{Christiano 2014 (protein turnover)}\newline Christiano R, Nagaraj N, Frohlich F, Walther TC. Cell Rep 2014;9:1959-1965.\newline \href{https://doi.org/10.1016/j.celrep.2014.10.065}{doi:10.1016/\allowbreak j.celrep.2014.10.065} & Scale band: tier 3, ordered by measurements (3,900). & Cell Reports SI (unconfirmed) \\
\addlinespace[5pt]
135 & \textbf{Kitzman 2015 (GAL4 DMS)}\newline Kitzman JO, Starita LM, Lo RS, Fields S, Shendure J. Nat Methods 2015;12:203-206.\newline \href{https://doi.org/10.1038/nmeth.3223}{doi:10.1038/\allowbreak nmeth.3223} & Scale band: tier 3, ordered by measurements (2,400). & Nat Methods SI (unconfirmed) \\
\addlinespace[5pt]
136 & \textbf{Sharon 2012 (designed promoter library)}\newline Sharon E, Kalma Y, Sharp A, et al., Segal E. Nat Biotechnol 2012;30:521-530.\newline \href{https://doi.org/10.1038/nbt.2205}{doi:10.1038/\allowbreak nbt.2205} & Scale band: tier 3, ordered by measurements (2,000). & Nat Biotechnol SI (unconfirmed) \\
\addlinespace[5pt]
137 & \textbf{Sadhu 2018 (CRISPR-directed mitotic recombination)}\newline Sadhu MJ, Bloom JS, Day L, Siegel JJ, Kosuri S, Kruglyak L. eLife 2018;7:e41090.\newline \href{https://doi.org/10.7554/eLife.41090}{doi:10.7554/\allowbreak eLife.41090} & Scale band: tier 3, ordered by measurements (1,500). & eLife data availability (unconfirmed) \\
\addlinespace[5pt]
138 & \textbf{Duveau 2021 (TDH3 promoter allele series)}\newline Duveau F, Vande Zande P, Metzger BP, et al., Wittkopp PJ. eLife 2021;10:e67806.\newline \href{https://doi.org/10.7554/eLife.67806}{doi:10.7554/\allowbreak eLife.67806} & Scale band: tier 3, ordered by measurements (1,000). & eLife data availability (unconfirmed) \\
\addlinespace[5pt]
139 & \textbf{Strope 2015 (100 clinical and natural genomes)}\newline Strope PK, Skelly DA, Kozmin SG, Mahadevan G, Stone EA, Magwene PM, Dietrich FS, McCusker JH. Genome Res 2015;25:762-774.\newline \href{https://doi.org/10.1101/gr.185538.114}{doi:10.1101/\allowbreak gr.185538.114} & Scale band: tier 3, ordered by measurements (1,000). & Genome Research SI; SRA (unconfirmed) \\
\addlinespace[5pt]
140 & \textbf{Schulte 2023 (mitochondrial complexome)}\newline Schulte U, et al. Nature 2023;614:153-159.\newline \href{https://doi.org/10.1038/s41586-022-05641-w}{doi:10.1038/\allowbreak s41586-022-05641-w} & Scale band: tier 3, ordered by measurements (900). & complexomics.org + ProteomeXchange \\
\addlinespace[5pt]
141 & \textbf{Metzger 2015 (promoter mutation effects on expression and noise)}\newline Metzger BPH, Yuan DC, Gruber JD, Duveau F, Wittkopp PJ. Nature 2015;521:344-346.\newline \href{https://doi.org/10.1038/nature14424}{doi:10.1038/\allowbreak nature14424} & Scale band: tier 3, ordered by measurements (470). & Nature SI (unconfirmed) \\
\addlinespace[5pt]
142 & \textbf{Venkataram 2016 (barcoded adaptive lineages, sequenced)}\newline Venkataram S, Dunn B, Li Y, et al., Petrov DA, Kryazhimskiy S, Sherlock G. Cell 2016;166:1585-1596.\newline \href{https://doi.org/10.1016/j.cell.2016.08.002}{doi:10.1016/\allowbreak j.cell.2016.08.002} & Scale band: tier 3, ordered by measurements (400). & Cell SI; SRA (unconfirmed) \\
\addlinespace[5pt]
143 & \textbf{Jakociunas 2021 (degron-tuned terpene flux)}\newline Jakociunas T, Klitgaard AK, Kristensen M, Jensen MK, Keasling JD. 2021.\newline \href{https://doi.org/10.1002/bit.27735}{doi:10.1002/\allowbreak bit.27735} & Scale band: tier 3, ordered by measurements (270). & journal SI (unconfirmed) \\
\addlinespace[5pt]
144 & \textbf{Shen 2016 (SCRaMbLE genome minimization)}\newline Shen Y, Stracquadanio G, Wang Y, et al., Boeke JD, Bader JS. Genome Res 2016;26:36-49.\newline \href{https://doi.org/10.1101/gr.193433.115}{doi:10.1101/\allowbreak gr.193433.115} & Scale band: tier 3, ordered by measurements (200). & Genome Research SI (unconfirmed) \\
\addlinespace[5pt]
145 & \textbf{Blount 2018 (SCRaMbLE for phenotype improvement)}\newline Blount BA, Gowers GF, Ho JCH, Ledesma-Amaro R, Jovicevic D, McKiernan RM, Xie ZX, Li BZ, Yuan YJ, Ellis T. Nat Commun 2018;9:1932.\newline \href{https://doi.org/10.1038/s41467-018-03143-w}{doi:10.1038/\allowbreak s41467-018-03143-w} & Scale band: tier 3, ordered by measurements (200). & Nat Commun SI (unconfirmed) \\
\addlinespace[5pt]
146 & \textbf{Liti 2009 (population genomics of S. cerevisiae and S. paradoxus)}\newline Liti G, Carter DM, Moses AM, et al., Louis EJ. Nature 2009;458:337-341.\newline \href{https://doi.org/10.1038/nature07743}{doi:10.1038/\allowbreak nature07743} & Scale band: tier 3, ordered by measurements (70). & SGRP /\allowbreak  ENA (unconfirmed) \\
\addlinespace[5pt]
147 & \textbf{Yue 2017 (wild vs domesticated assemblies)}\newline Yue JX, Li J, Aigrain L, et al., Liti G. Nat Genet 2017;49:913-924.\newline \href{https://doi.org/10.1038/ng.3847}{doi:10.1038/\allowbreak ng.3847} & Scale band: tier 3, ordered by measurements (12). & ENA/\allowbreak GenBank (accession unconfirmed) \\
\addlinespace[5pt]
148 & \textbf{Zackrisson 2016 (Scan-o-matic growth curves)}\newline Zackrisson M, Hallin J, Ottosson LG, et al., Warringer J, Blomberg A. G3 2016;6:3003-3014.\newline \href{https://doi.org/10.1534/g3.116.032342}{doi:10.1534/\allowbreak g3.116.032342} & Scale band: tier 3. & Open-source platform; per-study datasets \\
\addlinespace[5pt]
149 & \textbf{Spellman 1998 (cell-cycle transcriptome)}\newline Spellman PT, Sherlock G, Zhang MQ, Iyer VR, Anders K, Eisen MB, Brown PO, Botstein D, Futcher B. Mol Biol Cell 1998;9:3273-3297.\newline \href{https://doi.org/10.1091/mbc.9.12.3273}{doi:10.1091/\allowbreak mbc.9.12.3273} & Scale band: tier 4, ordered by measurements (462,000). & SGD /\allowbreak  SPELL \\
\addlinespace[5pt]
150 & \textbf{Levy 2008 (morphology and phenotypic capacitance)}\newline Levy SF, Siegal ML. PLoS Biol 2008;6:e264.\newline \href{https://doi.org/10.1371/journal.pbio.0060264}{doi:10.1371/\allowbreak journal.pbio.0060264} & Scale band: tier 4, ordered by measurements (340,000). & PLoS Biology SI (unconfirmed) \\
\addlinespace[5pt]
151 & \textbf{Causton 2001 (environmental change transcriptome)}\newline Causton HC, Ren B, Koh SS, et al., Young RA. Mol Biol Cell 2001;12:323-337.\newline \href{https://doi.org/10.1091/mbc.12.2.323}{doi:10.1091/\allowbreak mbc.12.2.323} & Scale band: tier 4, ordered by measurements (240,000). & SGD /\allowbreak  SPELL \\
\addlinespace[5pt]
152 & \textbf{Gasch 2001 (DNA damage response transcriptome)}\newline Gasch AP, Huang M, Metzner S, Botstein D, Elledge SJ, Brown PO. Mol Biol Cell 2001;12:2987-3003.\newline \href{https://doi.org/10.1091/mbc.12.10.2987}{doi:10.1091/\allowbreak mbc.12.10.2987} & Scale band: tier 4, ordered by measurements (144,000). & SGD /\allowbreak  SPELL \\
\addlinespace[5pt]
153 & \textbf{Handfield 2013 (image-based localization change)}\newline Handfield LF, Chong YT, Simmons J, Andrews BJ, Moses AM. PLoS Comput Biol 2013;9:e1003085.\newline \href{https://doi.org/10.1371/journal.pcbi.1003085}{doi:10.1371/\allowbreak journal.pcbi.1003085} & Scale band: tier 4, ordered by measurements (136,000). & PLoS Comput Biol SI (unconfirmed) \\
\addlinespace[5pt]
154 & \textbf{Kaplan 2009 (nucleosome sequence preferences)}\newline Kaplan N, Moore IK, Fondufe-Mittendorf Y, et al., Segal E. Nature 2009;458:362-366.\newline \href{https://doi.org/10.1038/nature07667}{doi:10.1038/\allowbreak nature07667} & Scale band: tier 4, ordered by measurements (134,000). & GEO (unconfirmed) \\
\addlinespace[5pt]
155 & \textbf{Giaever 2004 (chemogenomic drug-target profiling)}\newline Giaever G, Flaherty P, Kumm J, et al., Davis RW. PNAS 2004;101:793-798.\newline \href{https://doi.org/10.1073/pnas.0307490100}{doi:10.1073/\allowbreak pnas.0307490100} & Scale band: tier 4, ordered by measurements (59,000). & PNAS SI (unconfirmed) \\
\addlinespace[5pt]
156 & \textbf{Paulo 2016 (TMT proteome across carbon sources)}\newline Paulo JA, O'Connell JD, Gaun A, Gygi SP. Mol Biol Cell 2016;27:2823-2838.\newline \href{https://doi.org/10.1091/mbc.E16-03-0184}{doi:10.1091/\allowbreak mbc.E16-03-0184} & Scale band: tier 4, ordered by measurements (22,500). & ProteomeXchange (unconfirmed) \\
\addlinespace[5pt]
157 & \textbf{Gibney 2013 (osmotic and glucose stress fitness)}\newline Gibney PA, Lu C, Caudy AA, Hess DC, Botstein D. PNAS 2013;110:E4393-E4402.\newline \href{https://doi.org/10.1073/pnas.1318100110}{doi:10.1073/\allowbreak pnas.1318100110} & Scale band: tier 4, ordered by measurements (19,200). & PNAS SI (unconfirmed) \\
\addlinespace[5pt]
158 & \textbf{Meurer 2018 (genome-wide GFP library, rebuilt)}\newline Meurer M, Duan Y, Sass E, Kats I, Herbst K, Buchmuller BC, Dederer V, Huber F, Kirrmaier D, Stefl M, Van Laer K, Dick TP, Lemberg MK, Khmelinskii A, Levy ED, Knop M. Nat Methods 2018;15:617-620.\newline \href{https://doi.org/10.1038/s41592-018-0044-9}{doi:10.1038/\allowbreak s41592-018-0044-9} & Scale band: tier 4, ordered by measurements (12,000). & Nat Methods SI (DOI collision with Weill 2018 UNRESOLVED) \\
\addlinespace[5pt]
159 & \textbf{Bennett 2001 (ionizing-radiation resistance)}\newline Bennett CB, et al. Nat Genet 2001;29:426-434.\newline \href{https://doi.org/10.1038/ng778}{doi:10.1038/\allowbreak ng778} & Scale band: tier 4, ordered by measurements (9,600). & Nat Genet SI (likely PDF) \\
\addlinespace[5pt]
160 & \textbf{Sugiyama (lactic-acid tolerance)}\newline Sugiyama M, Kaneko Y, et al. National Research Institute of Brewing technical report.\newline \href{https://www.nisr.or.jp/wp-content/uploads/NISR06sugiyama.pdf}{nisr.or.jp/\allowbreak wp-content/\allowbreak uploads/\allowbreak NISR06sugiyama.pdf} & Scale band: tier 4, ordered by measurements (9,600). & NRIB technical report PDF \\
\addlinespace[5pt]
161 & \textbf{Sousa 2013 (acetic-acid programmed cell death screen)}\newline Sousa M, Duarte AM, Fernandes TR, et al. 2013. PMID 24286259.\newline \href{https://pubmed.ncbi.nlm.nih.gov/24286259/}{PMID 24286259} & Scale band: tier 4, ordered by measurements (9,600). & journal SI (unconfirmed) \\
\addlinespace[5pt]
162 & \textbf{Nagalakshmi 2008 (RNA-seq transcriptome annotation)}\newline Nagalakshmi U, Wang Z, Waern K, Shou C, Raha D, Gerstein M, Snyder M. Science 2008;320:1344-1349.\newline \href{https://doi.org/10.1126/science.1158441}{doi:10.1126/\allowbreak science.1158441} & Scale band: tier 4, ordered by measurements (6,000). & SRA /\allowbreak  GEO (unconfirmed) \\
\addlinespace[5pt]
163 & \textbf{Yofe 2016 (SWAT N-terminal library)}\newline Yofe I, Weill U, Meurer M, et al., Knop M, Schuldiner M. Nat Methods 2016;13:371-378.\newline \href{https://doi.org/10.1038/nmeth.3795}{doi:10.1038/\allowbreak nmeth.3795} & Scale band: tier 4, ordered by measurements (5,400). & Nat Methods SI; EUROSCARF (unconfirmed) \\
\addlinespace[5pt]
164 & \textbf{Hénault 2023 (hybrid and introgression panel)}\newline Henault M, Marsit S, Charron G, Landry CR. Nat Commun / Landry lab hybrid panels.\newline \href{https://doi.org/10.1038/s41467-024-46155-5}{doi:10.1038/\allowbreak s41467-024-46155-5} & Scale band: tier 4, ordered by measurements (5,000). & journal SI (unconfirmed) \\
\addlinespace[5pt]
165 & \textbf{Cubillos 2013 (F12 recombinant panel construction)}\newline Cubillos FA, Parts L, Salinas F, et al., Liti G. Genetics 2013;195:1141-1155.\newline \href{https://doi.org/10.1534/genetics.113.155515}{doi:10.1534/\allowbreak genetics.113.155515} & Scale band: tier 4, ordered by measurements (3,840). & Genetics SI Table S1 \\
\addlinespace[5pt]
166 & \textbf{Kessi-Perez 2019 (nitrogen-consumption QTL panel)}\newline Kessi-Perez EI, Molinet J, Martinez C, et al. 2019/2020 nitrogen-utilization QTL work.\newline \href{https://doi.org/10.1186/s13068-019-1440-9}{doi:10.1186/\allowbreak s13068-019-1440-9} & Scale band: tier 4, ordered by measurements (200). & journal SI (unconfirmed) \\
\addlinespace[5pt]
167 & \textbf{Lian 2016 (combinatorial xylose pathway optimization)}\newline Lian J, Chao R, Zhao H. Metab Eng 2016;33:139-148.\newline \href{https://doi.org/10.1016/j.ymben.2015.11.001}{doi:10.1016/\allowbreak j.ymben.2015.11.001} & Scale band: tier 4, ordered by measurements (100). & Metab Eng SI (unconfirmed) \\
\addlinespace[5pt]
\end{longtable}
\endgroup
\end{landscape}
